%% ============================================================
%% koushin-history.tex --- 厚胤塾 歴史（日本史・世界史）共通プリアンブル
%% 使い方: 通常はこのファイルを直接 \input せず、
%%   \input{koushin-nihonshi}  または  %% ============================================================
%% koushin-sekaishi.tex --- 厚胤塾 世界史用プリアンブル
%% 使い方: %% ============================================================
%% koushin-base.tex --- 厚胤塾 共通プリアンブル（全科目共通）
%% ------------------------------------------------------------
%% 対応環境: platex + dvipdfmx（Cloud LaTeX 標準構成）
%%
%% 使い方:
%%   \documentclass[dvipdfmx]{jsbook}   % jsarticle でも可
%%   %% ============================================================
%% koushin-base.tex --- 厚胤塾 共通プリアンブル（全科目共通）
%% ------------------------------------------------------------
%% 対応環境: platex + dvipdfmx（Cloud LaTeX 標準構成）
%%
%% 使い方:
%%   \documentclass[dvipdfmx]{jsbook}   % jsarticle でも可
%%   \input{koushin-base}      % ← 必ず最初に読み込む
%%   \input{koushin-physics}   % ← 科目別プリアンブル（1つだけ）
%%   \begin{document}
%%   ...
%%
%% ブランドカラー: mainBlue = RGB(0,51,102)  ← 全科目共通・不変
%% 科目色       : accentColor ← 科目別ファイルが上書きする
%% ============================================================

%% ---------- 基本パッケージ ----------
\usepackage{graphicx}
\usepackage{wrapfig}
\usepackage{xcolor}
\usepackage{amsmath,amsthm,amssymb}
\usepackage{bm}
\usepackage{enumitem}
\usepackage{fancybox}
\usepackage{multicol}
\usepackage{pdfpages}
\usepackage{float}
\usepackage{marginnote}
\usepackage{ascmac}
\usepackage{url}
\usepackage{tikz}
\usetikzlibrary{patterns}
\usepackage[most]{tcolorbox}
\tcbuselibrary{skins,breakable,xparse}
\usepackage[normalem]{ulem}
\usepackage{tocloft}
\usepackage{fancyhdr}

%% ---------- ブランドカラー ----------
\definecolor{mainBlue}{RGB}{0,51,102}      % 厚胤塾ブランド色（固定）
\definecolor{accentColor}{RGB}{0,51,102}   % 科目色（科目別ファイルで上書き）
\providecommand{\SubjectName}{}            % 科目名（科目別ファイルで上書き）

%% ---------- 余白 ----------
\usepackage[top=25mm,bottom=25mm,left=25mm,right=25mm]{geometry}

%% ---------- 2段組の設定 ----------
\setlength{\columnsep}{6mm}
\setlength{\columnseprule}{0.4pt}
\def\columnseprulecolor{\color{mainBlue}}

%% ---------- 丸数字 \maru{}（platexでUnicode①が使えないため） ----------
\DeclareRobustCommand{\maru}[1]{%
  \tikz[baseline=(char.base)]{%
    \node[shape=circle,draw,inner sep=0.6pt,line width=0.4pt] (char)
      {\footnotesize #1};}}

%% ---------- 式番号・図表番号（丸数字＆通し番号） ----------
\makeatletter
\@ifundefined{c@chapter}{}{%
  \counterwithout{equation}{chapter}%
  \counterwithout{figure}{chapter}%
  \counterwithout{table}{chapter}%
}
\makeatother
\renewcommand{\theequation}{\maru{\arabic{equation}}}
\renewcommand{\thefigure}{\arabic{figure}}
\renewcommand{\thetable}{\arabic{table}}

%% ---------- enumerate のラベル形式 ----------
\setlist[enumerate,1]{label=(\arabic*),  font=\normalfont\upshape}
\setlist[enumerate,2]{label=(\roman*).,  font=\normalfont\upshape}
\setlist[enumerate,3]{label=(\alph*),    font=\normalfont\upshape}

%% ---------- 目次デザイン ----------
\renewcommand{\contentsname}{内容一覧}
\setcounter{tocdepth}{2}
\renewcommand{\cftpartfont}{\LARGE\bfseries\color{mainBlue}}
\renewcommand{\cftpartpagefont}{\LARGE\bfseries\color{mainBlue}}
\renewcommand{\cftsecfont}{\Large\bfseries\color{accentColor}}
\renewcommand{\cftsecpagefont}{\Large\bfseries\color{accentColor}}
\setlength{\cftbeforesecskip}{6pt}

%% ---------- ヘッダー・フッター ----------
\pagestyle{fancy}
\fancyhf{}
\fancyhead[LE]{\textcolor{accentColor}{\textsf{\leftmark}}}
\fancyhead[RO]{\textcolor{accentColor}{\textsf{\rightmark}}}
\fancyfoot[C]{\textcolor{mainBlue}{\thepage}}
\fancyfoot[LE,RO]{\textcolor{accentColor}{\scriptsize 厚胤塾\ \SubjectName}}
\renewcommand{\headrulewidth}{0pt}
\renewcommand{\footrulewidth}{0pt}
\renewcommand{\headrule}{%
  \hbox to \headwidth{\color{accentColor}\leaders\hrule height 0.5pt\hfill}}
\renewcommand{\footrule}{%
  \hbox to \headwidth{\color{accentColor}\leaders\hrule height 0.5pt\hfill}}

%% ============================================================
%%  装飾枠（tcolorbox）--- 枠色は accentColor に連動
%% ============================================================

%% 授業例・例示用（白背景・影つき・タイトルあり）
\newtcolorbox{jyurei}[1][]{
  enhanced, colframe=accentColor, colback=white,
  coltitle=white, fonttitle=\bfseries, drop fuzzy shadow,
  title=#1, boxrule=1pt, arc=4pt, breakable
}

%% グレー枠・タイトル浮き出し
\newtcolorbox{mybox}[1][]{
  enhanced, colframe=white!35!black, colback=white,
  title=#1, fonttitle=\bfseries, breakable=true,
  coltitle=white!20!black,
  attach boxed title to top left={xshift=5mm, yshift=-3mm},
  boxed title style={colframe=white!35!black, colback=white},
  top=4mm
}

%% 科目色枠・クリーム背景・タイトルボックス浮き出し
\newtcolorbox{b_s_b_t}[1][]{
  colback=yellow!10, colframe=accentColor, coltitle=white,
  fonttitle=\bfseries, title=#1, boxrule=1pt, arc=4pt,
  breakable, enhanced,
  attach boxed title to top left={xshift=1em,
    yshift*=-\tcboxedtitleheight/2},
  colbacktitle=accentColor,
  boxed title style={frame code={%
    \path[fill=accentColor] (frame.south west)--(frame.north west)--
    (frame.north east)--(frame.south east)--cycle;},
    interior hidden},
}

%% 科目色実線枠・クリーム背景・シンプルタイトル
\newtcolorbox{bluebox}[1][]{
  colframe=accentColor, colback=yellow!10, coltitle=white,
  fonttitle=\bfseries, title=#1, boxrule=1pt, arc=4pt, breakable
}

%% 破線枠・クリーム背景・タイトルあり
\newenvironment{mybox_d_b}[1]{%
  \begin{tcolorbox}[
    coltitle=black, title=#1, colback=yellow!10,
    frame hidden, interior hidden, enhanced, arc=1em,
    borderline={1pt}{0pt}{dashed},
    attach boxed title to top left={xshift=1em,
      yshift*=-\tcboxedtitleheight/2},
    colbacktitle=white, boxed title style={frame hidden},
    breakable,
  ]
}{\end{tcolorbox}}

%% 破線枠・タイトルなし・背景なし
\newenvironment{mybox_dash}{%
  \begin{tcolorbox}[
    frame hidden, interior hidden, enhanced, arc=1em,
    borderline={1pt}{0pt}{dashed}
  ]
}{\end{tcolorbox}}

%% 破線枠・タイトルあり・背景なし（解法など）
\newenvironment{kaiho}[1]{%
  \begin{tcolorbox}[
    coltitle=black, title=#1, frame hidden, interior hidden,
    enhanced, arc=1em, borderline={1pt}{0pt}{dashed},
    attach boxed title to top left={xshift=1em,
      yshift*=-\tcboxedtitleheight/2},
    colbacktitle=white, boxed title style={frame hidden},
  ]
}{\end{tcolorbox}}

%% 黒実線・タイトル切り欠き枠（原理・法則）
\DeclareTColorBox{genri}{o m O{.5} O{}}{%
  empty, left=4mm, right=4mm, top=-1mm,
  attach boxed title to top left={xshift=1.2zw},
  boxed title style={empty, left=-2mm, right=-2mm},
  colframe=black, coltitle=black, coltext=black, breakable,
  underlay unbroken={%
    \draw[black, line width=#3pt]
      (title.east) -- (title.east-|frame.east) --
      (frame.south east) -- (frame.south west) --
      (title.west-|frame.west) -- (title.west);},
  underlay first={%
    \draw[black, line width=#3pt]
      (title.east) -- (title.east-|frame.east) -- (frame.south east);
    \draw[black, line width=#3pt]
      (frame.south west) -- (title.west-|frame.west) -- (title.west);},
  underlay middle={%
    \draw[black, line width=#3pt] (frame.north east) -- (frame.south east);
    \draw[black, line width=#3pt] (frame.south west) -- (frame.north west);},
  underlay last={%
    \draw[black, line width=#3pt]
      (frame.north east) -- (frame.south east) --
      (frame.south west) -- (frame.north west);},
  fonttitle=\gtfamily,
  IfValueTF={#1}{title=【#2】〈#1〉}{title=【#2】},
  #4
}

%% 警告枠（赤：全科目共通で赤固定）
\newtcolorbox{keikoku}[1][]{
  colframe=red!70!black, colback=yellow!10, coltitle=white,
  fonttitle=\bfseries, title=#1, boxrule=1pt, arc=4pt,
  breakable, enhanced,
  attach boxed title to top left={xshift=1em,
    yshift*=-\tcboxedtitleheight/2},
  colbacktitle=red!70!black,
  boxed title style={frame code={%
    \path[fill=red!70!black] (frame.south west)--(frame.north west)--
    (frame.north east)--(frame.south east)--cycle;},
    interior hidden},
}

%% 要点まとめ枠（科目色・薄い科目色背景）
\newtcolorbox{pointbox}[1][ポイント]{
  enhanced, colframe=accentColor, colback=accentColor!4!white,
  coltitle=white, colbacktitle=accentColor, fonttitle=\bfseries,
  title=#1, boxrule=1pt, arc=3pt, breakable,
  attach boxed title to top left={xshift=1em,
    yshift*=-\tcboxedtitleheight/2}
}

%% ============================================================
%%  見出しデザイン
%% ============================================================

%% --- セクション（番号白抜き＋帯） ---
\makeatletter
\renewcommand{\thesection}{\arabic{section}}
\newcommand\Section[1]{%
  \Section@topmargin%
  \refstepcounter{section}%
  {\setlength{\parindent}{0pt}%
   \Section@font%
   \Section@title{\thesection}{#1}%
   \par\nobreak}%
  \Section@bottommargin%
}
\newcommand{\Section@topmargin}{\vspace{2\Cvs}}
\newcommand{\Section@bottommargin}{\vspace{1\Cvs}}
\newcommand{\Section@font}{\headfont\Large}
\newcommand{\Section@title}[2]{%
  \textcolor{accentColor!30}{\rlap{\rule[-0.6zh]{\textwidth}{2zh}}}%
  \textcolor{accentColor}{\rlap{\rule[-0.6zh]{3.5zw}{2zh}}}%
  \makebox[3.5zw][c]{\textcolor{white}{#1}}%
  \hspace{0.9zw}#2%
}
\makeatother

%% --- 太バー付き見出し ---
\newcommand{\DecoratedSection}[1]{%
\vspace{2\Cvs}%
{\setlength{\parindent}{0pt}%
\begin{tcolorbox}[
  enhanced, colframe=accentColor, colback=white,
  boxrule=1pt, arc=0mm, left=10pt, right=10pt, top=8pt, bottom=8pt,
  leftrule=12pt, width=\textwidth, fontupper=\headfont\Large
]{#1}\end{tcolorbox}}%
\vspace{\Cvs}}

%% --- 二重ライン見出し ---
\newcommand{\FancySection}[1]{%
\vspace{2\Cvs}%
{\setlength{\parindent}{0pt}%
\begin{tcolorbox}[
  enhanced, colframe=accentColor, colback=white,
  boxrule=1pt, arc=3mm, left=10pt, right=10pt, top=8pt, bottom=8pt,
  leftrule=0pt, width=\textwidth,
  overlay={%
    \draw[accentColor, line width=3pt]
      ([xshift=2pt]frame.north west) -- ([xshift=-2pt]frame.north east);
    \draw[accentColor!30, line width=1.5pt]
      ([xshift=2pt, yshift=-3pt]frame.north west) --
      ([xshift=-2pt, yshift=-3pt]frame.north east);},
  fontupper=\headfont\Large
]{#1}\end{tcolorbox}}%
\vspace{\Cvs}}

%% --- 入試問題バナー ---
\newcommand{\ChallengeTitle}{%
\vspace{0.3cm}
\begin{tcolorbox}[
  enhanced, colframe=accentColor, colback=white,
  boxrule=0pt, arc=0mm, left=20pt, right=20pt, top=8pt, bottom=8pt,
  width=\textwidth,
  overlay={%
    \fill[accentColor]
      ([xshift=-3pt]frame.north west) rectangle (frame.south west);
    \fill[accentColor!90]
      (frame.north west) rectangle ([xshift=\textwidth]frame.north east);
    \fill[left color=accentColor!20, right color=white]
      ([yshift=-3pt]frame.north west) rectangle
      ([xshift=8cm, yshift=-4pt]frame.north west);
    \node[anchor=west, inner sep=0pt, font=\Large\bfseries]
      at ([xshift=25pt, yshift=-14pt]frame.north west)
      {\textcolor{accentColor}{入試問題に挑戦}};},
  height=2cm
]\end{tcolorbox}
\vspace{0.3cm}}

%% --- テスト用タイトルバナー（#1: テスト名, #2: 範囲） ---
\newcommand{\KoushinTestTitle}[2]{%
\begin{tcolorbox}[
  enhanced, colframe=mainBlue, colback=white,
  boxrule=1pt, arc=0mm, left=16pt, right=12pt, top=10pt, bottom=10pt,
  overlay={%
    \fill[accentColor]
      (frame.north west) rectangle ([xshift=6pt]frame.south west);}
]
{\headfont\LARGE\textcolor{mainBlue}{#1}}\hfill
{\headfont\large\textcolor{accentColor}{\SubjectName}}\par
\vspace{2pt}
{\large #2}\hfill{\small\textcolor{mainBlue}{厚胤塾}}\par
\vspace{8pt}
\hfill 氏名\ \uline{\hspace{12zw}}\hspace{1zw}%
得点\ \fbox{\rule{0pt}{1.6zh}\hspace{5zw}}
\end{tcolorbox}
\vspace{0.5\Cvs}}

%% ============================================================
%%  定理環境（問題・解答系）
%% ============================================================
\newtheorem{problem}{問題}
\newtheorem{reidai}{例題}
\newtheorem{ans}{答え}
\newtheorem{ensyu}{演習}
\newtheorem{solution}{解答}
\newtheorem{explanation}{解説}
\newtheorem{kakunin}{確認問題}
\newtheorem{kakusolu}{解説}

%% ============================================================
%%  アンダーライン
%% ============================================================
\font\sixly=lasy6
\newcommand{\waveunderline}[1]{%
  \bgroup\markoverwith{\lower3.5pt\hbox{\sixly \char58}}\ULon{#1}}
\newcommand{\dashunderline}[1]{%
  \bgroup\markoverwith{\hbox to .5em{\hrulefill}}\ULon{#1}}
\newcommand{\solidunderline}[1]{\uline{#1}}

%% ============================================================
%%  穴埋め \anaume 系
%% ============================================================
\setlength{\fboxsep}{1pt}
%% 太枠（初出の解答欄）
\newcommand{\anaume}[1]{\setlength{\fboxrule}{1pt}%
  \framebox[40pt][c]{\mbox{\textbf{\large{#1}}}}\setlength{\fboxrule}{0.4pt}}
%% 細枠（2回目以降）
\newcommand{\sanaume}[1]{\framebox[40pt][c]{\mbox{\large{#1}}}}
%% 横長太枠（3桁以上）
\newcommand{\banaume}[1]{\setlength{\fboxrule}{1pt}%
  \fbox{\mbox{\textbf{\large{\phantom{a}#1\phantom{a}}}}}\setlength{\fboxrule}{0.4pt}}
%% 横長細枠
\newcommand{\sbanaume}[1]{\fbox{\mbox{\large{\phantom{a}#1\phantom{a}}}}}
%% 指数用
\newcommand{\sisuuanaume}[1]{\setlength{\fboxrule}{1pt}%
  \framebox[16pt][c]{\mbox{\textbf{#1}}}\setlength{\fboxrule}{0.4pt}}
\newcommand{\ssisuuanaume}[1]{\framebox[16pt][c]{\mbox{#1}}}
%% 選択肢用（二重枠）
\newcommand{\senntakuanaume}[1]{\setlength{\fboxrule}{0.6pt}%
  \doublebox{\mbox{\textbf{\phantom{ア}#1\phantom{ア}}}}\setlength{\fboxrule}{0.4pt}}
\newcommand{\ssenntakuanaume}[1]{\doublebox{\mbox{\phantom{ア}#1\phantom{ア}}}}

%% ============================================================
%%  hyperref（最後に読み込む）
%% ============================================================
\usepackage{hyperref}
\usepackage{pxjahyper}
\hypersetup{hidelinks}
      % ← 必ず最初に読み込む
%%   \input{koushin-physics}   % ← 科目別プリアンブル（1つだけ）
%%   \begin{document}
%%   ...
%%
%% ブランドカラー: mainBlue = RGB(0,51,102)  ← 全科目共通・不変
%% 科目色       : accentColor ← 科目別ファイルが上書きする
%% ============================================================

%% ---------- 基本パッケージ ----------
\usepackage{graphicx}
\usepackage{wrapfig}
\usepackage{xcolor}
\usepackage{amsmath,amsthm,amssymb}
\usepackage{bm}
\usepackage{enumitem}
\usepackage{fancybox}
\usepackage{multicol}
\usepackage{pdfpages}
\usepackage{float}
\usepackage{marginnote}
\usepackage{ascmac}
\usepackage{url}
\usepackage{tikz}
\usetikzlibrary{patterns}
\usepackage[most]{tcolorbox}
\tcbuselibrary{skins,breakable,xparse}
\usepackage[normalem]{ulem}
\usepackage{tocloft}
\usepackage{fancyhdr}

%% ---------- ブランドカラー ----------
\definecolor{mainBlue}{RGB}{0,51,102}      % 厚胤塾ブランド色（固定）
\definecolor{accentColor}{RGB}{0,51,102}   % 科目色（科目別ファイルで上書き）
\providecommand{\SubjectName}{}            % 科目名（科目別ファイルで上書き）

%% ---------- 余白 ----------
\usepackage[top=25mm,bottom=25mm,left=25mm,right=25mm]{geometry}

%% ---------- 2段組の設定 ----------
\setlength{\columnsep}{6mm}
\setlength{\columnseprule}{0.4pt}
\def\columnseprulecolor{\color{mainBlue}}

%% ---------- 丸数字 \maru{}（platexでUnicode①が使えないため） ----------
\DeclareRobustCommand{\maru}[1]{%
  \tikz[baseline=(char.base)]{%
    \node[shape=circle,draw,inner sep=0.6pt,line width=0.4pt] (char)
      {\footnotesize #1};}}

%% ---------- 式番号・図表番号（丸数字＆通し番号） ----------
\makeatletter
\@ifundefined{c@chapter}{}{%
  \counterwithout{equation}{chapter}%
  \counterwithout{figure}{chapter}%
  \counterwithout{table}{chapter}%
}
\makeatother
\renewcommand{\theequation}{\maru{\arabic{equation}}}
\renewcommand{\thefigure}{\arabic{figure}}
\renewcommand{\thetable}{\arabic{table}}

%% ---------- enumerate のラベル形式 ----------
\setlist[enumerate,1]{label=(\arabic*),  font=\normalfont\upshape}
\setlist[enumerate,2]{label=(\roman*).,  font=\normalfont\upshape}
\setlist[enumerate,3]{label=(\alph*),    font=\normalfont\upshape}

%% ---------- 目次デザイン ----------
\renewcommand{\contentsname}{内容一覧}
\setcounter{tocdepth}{2}
\renewcommand{\cftpartfont}{\LARGE\bfseries\color{mainBlue}}
\renewcommand{\cftpartpagefont}{\LARGE\bfseries\color{mainBlue}}
\renewcommand{\cftsecfont}{\Large\bfseries\color{accentColor}}
\renewcommand{\cftsecpagefont}{\Large\bfseries\color{accentColor}}
\setlength{\cftbeforesecskip}{6pt}

%% ---------- ヘッダー・フッター ----------
\pagestyle{fancy}
\fancyhf{}
\fancyhead[LE]{\textcolor{accentColor}{\textsf{\leftmark}}}
\fancyhead[RO]{\textcolor{accentColor}{\textsf{\rightmark}}}
\fancyfoot[C]{\textcolor{mainBlue}{\thepage}}
\fancyfoot[LE,RO]{\textcolor{accentColor}{\scriptsize 厚胤塾\ \SubjectName}}
\renewcommand{\headrulewidth}{0pt}
\renewcommand{\footrulewidth}{0pt}
\renewcommand{\headrule}{%
  \hbox to \headwidth{\color{accentColor}\leaders\hrule height 0.5pt\hfill}}
\renewcommand{\footrule}{%
  \hbox to \headwidth{\color{accentColor}\leaders\hrule height 0.5pt\hfill}}

%% ============================================================
%%  装飾枠（tcolorbox）--- 枠色は accentColor に連動
%% ============================================================

%% 授業例・例示用（白背景・影つき・タイトルあり）
\newtcolorbox{jyurei}[1][]{
  enhanced, colframe=accentColor, colback=white,
  coltitle=white, fonttitle=\bfseries, drop fuzzy shadow,
  title=#1, boxrule=1pt, arc=4pt, breakable
}

%% グレー枠・タイトル浮き出し
\newtcolorbox{mybox}[1][]{
  enhanced, colframe=white!35!black, colback=white,
  title=#1, fonttitle=\bfseries, breakable=true,
  coltitle=white!20!black,
  attach boxed title to top left={xshift=5mm, yshift=-3mm},
  boxed title style={colframe=white!35!black, colback=white},
  top=4mm
}

%% 科目色枠・クリーム背景・タイトルボックス浮き出し
\newtcolorbox{b_s_b_t}[1][]{
  colback=yellow!10, colframe=accentColor, coltitle=white,
  fonttitle=\bfseries, title=#1, boxrule=1pt, arc=4pt,
  breakable, enhanced,
  attach boxed title to top left={xshift=1em,
    yshift*=-\tcboxedtitleheight/2},
  colbacktitle=accentColor,
  boxed title style={frame code={%
    \path[fill=accentColor] (frame.south west)--(frame.north west)--
    (frame.north east)--(frame.south east)--cycle;},
    interior hidden},
}

%% 科目色実線枠・クリーム背景・シンプルタイトル
\newtcolorbox{bluebox}[1][]{
  colframe=accentColor, colback=yellow!10, coltitle=white,
  fonttitle=\bfseries, title=#1, boxrule=1pt, arc=4pt, breakable
}

%% 破線枠・クリーム背景・タイトルあり
\newenvironment{mybox_d_b}[1]{%
  \begin{tcolorbox}[
    coltitle=black, title=#1, colback=yellow!10,
    frame hidden, interior hidden, enhanced, arc=1em,
    borderline={1pt}{0pt}{dashed},
    attach boxed title to top left={xshift=1em,
      yshift*=-\tcboxedtitleheight/2},
    colbacktitle=white, boxed title style={frame hidden},
    breakable,
  ]
}{\end{tcolorbox}}

%% 破線枠・タイトルなし・背景なし
\newenvironment{mybox_dash}{%
  \begin{tcolorbox}[
    frame hidden, interior hidden, enhanced, arc=1em,
    borderline={1pt}{0pt}{dashed}
  ]
}{\end{tcolorbox}}

%% 破線枠・タイトルあり・背景なし（解法など）
\newenvironment{kaiho}[1]{%
  \begin{tcolorbox}[
    coltitle=black, title=#1, frame hidden, interior hidden,
    enhanced, arc=1em, borderline={1pt}{0pt}{dashed},
    attach boxed title to top left={xshift=1em,
      yshift*=-\tcboxedtitleheight/2},
    colbacktitle=white, boxed title style={frame hidden},
  ]
}{\end{tcolorbox}}

%% 黒実線・タイトル切り欠き枠（原理・法則）
\DeclareTColorBox{genri}{o m O{.5} O{}}{%
  empty, left=4mm, right=4mm, top=-1mm,
  attach boxed title to top left={xshift=1.2zw},
  boxed title style={empty, left=-2mm, right=-2mm},
  colframe=black, coltitle=black, coltext=black, breakable,
  underlay unbroken={%
    \draw[black, line width=#3pt]
      (title.east) -- (title.east-|frame.east) --
      (frame.south east) -- (frame.south west) --
      (title.west-|frame.west) -- (title.west);},
  underlay first={%
    \draw[black, line width=#3pt]
      (title.east) -- (title.east-|frame.east) -- (frame.south east);
    \draw[black, line width=#3pt]
      (frame.south west) -- (title.west-|frame.west) -- (title.west);},
  underlay middle={%
    \draw[black, line width=#3pt] (frame.north east) -- (frame.south east);
    \draw[black, line width=#3pt] (frame.south west) -- (frame.north west);},
  underlay last={%
    \draw[black, line width=#3pt]
      (frame.north east) -- (frame.south east) --
      (frame.south west) -- (frame.north west);},
  fonttitle=\gtfamily,
  IfValueTF={#1}{title=【#2】〈#1〉}{title=【#2】},
  #4
}

%% 警告枠（赤：全科目共通で赤固定）
\newtcolorbox{keikoku}[1][]{
  colframe=red!70!black, colback=yellow!10, coltitle=white,
  fonttitle=\bfseries, title=#1, boxrule=1pt, arc=4pt,
  breakable, enhanced,
  attach boxed title to top left={xshift=1em,
    yshift*=-\tcboxedtitleheight/2},
  colbacktitle=red!70!black,
  boxed title style={frame code={%
    \path[fill=red!70!black] (frame.south west)--(frame.north west)--
    (frame.north east)--(frame.south east)--cycle;},
    interior hidden},
}

%% 要点まとめ枠（科目色・薄い科目色背景）
\newtcolorbox{pointbox}[1][ポイント]{
  enhanced, colframe=accentColor, colback=accentColor!4!white,
  coltitle=white, colbacktitle=accentColor, fonttitle=\bfseries,
  title=#1, boxrule=1pt, arc=3pt, breakable,
  attach boxed title to top left={xshift=1em,
    yshift*=-\tcboxedtitleheight/2}
}

%% ============================================================
%%  見出しデザイン
%% ============================================================

%% --- セクション（番号白抜き＋帯） ---
\makeatletter
\renewcommand{\thesection}{\arabic{section}}
\newcommand\Section[1]{%
  \Section@topmargin%
  \refstepcounter{section}%
  {\setlength{\parindent}{0pt}%
   \Section@font%
   \Section@title{\thesection}{#1}%
   \par\nobreak}%
  \Section@bottommargin%
}
\newcommand{\Section@topmargin}{\vspace{2\Cvs}}
\newcommand{\Section@bottommargin}{\vspace{1\Cvs}}
\newcommand{\Section@font}{\headfont\Large}
\newcommand{\Section@title}[2]{%
  \textcolor{accentColor!30}{\rlap{\rule[-0.6zh]{\textwidth}{2zh}}}%
  \textcolor{accentColor}{\rlap{\rule[-0.6zh]{3.5zw}{2zh}}}%
  \makebox[3.5zw][c]{\textcolor{white}{#1}}%
  \hspace{0.9zw}#2%
}
\makeatother

%% --- 太バー付き見出し ---
\newcommand{\DecoratedSection}[1]{%
\vspace{2\Cvs}%
{\setlength{\parindent}{0pt}%
\begin{tcolorbox}[
  enhanced, colframe=accentColor, colback=white,
  boxrule=1pt, arc=0mm, left=10pt, right=10pt, top=8pt, bottom=8pt,
  leftrule=12pt, width=\textwidth, fontupper=\headfont\Large
]{#1}\end{tcolorbox}}%
\vspace{\Cvs}}

%% --- 二重ライン見出し ---
\newcommand{\FancySection}[1]{%
\vspace{2\Cvs}%
{\setlength{\parindent}{0pt}%
\begin{tcolorbox}[
  enhanced, colframe=accentColor, colback=white,
  boxrule=1pt, arc=3mm, left=10pt, right=10pt, top=8pt, bottom=8pt,
  leftrule=0pt, width=\textwidth,
  overlay={%
    \draw[accentColor, line width=3pt]
      ([xshift=2pt]frame.north west) -- ([xshift=-2pt]frame.north east);
    \draw[accentColor!30, line width=1.5pt]
      ([xshift=2pt, yshift=-3pt]frame.north west) --
      ([xshift=-2pt, yshift=-3pt]frame.north east);},
  fontupper=\headfont\Large
]{#1}\end{tcolorbox}}%
\vspace{\Cvs}}

%% --- 入試問題バナー ---
\newcommand{\ChallengeTitle}{%
\vspace{0.3cm}
\begin{tcolorbox}[
  enhanced, colframe=accentColor, colback=white,
  boxrule=0pt, arc=0mm, left=20pt, right=20pt, top=8pt, bottom=8pt,
  width=\textwidth,
  overlay={%
    \fill[accentColor]
      ([xshift=-3pt]frame.north west) rectangle (frame.south west);
    \fill[accentColor!90]
      (frame.north west) rectangle ([xshift=\textwidth]frame.north east);
    \fill[left color=accentColor!20, right color=white]
      ([yshift=-3pt]frame.north west) rectangle
      ([xshift=8cm, yshift=-4pt]frame.north west);
    \node[anchor=west, inner sep=0pt, font=\Large\bfseries]
      at ([xshift=25pt, yshift=-14pt]frame.north west)
      {\textcolor{accentColor}{入試問題に挑戦}};},
  height=2cm
]\end{tcolorbox}
\vspace{0.3cm}}

%% --- テスト用タイトルバナー（#1: テスト名, #2: 範囲） ---
\newcommand{\KoushinTestTitle}[2]{%
\begin{tcolorbox}[
  enhanced, colframe=mainBlue, colback=white,
  boxrule=1pt, arc=0mm, left=16pt, right=12pt, top=10pt, bottom=10pt,
  overlay={%
    \fill[accentColor]
      (frame.north west) rectangle ([xshift=6pt]frame.south west);}
]
{\headfont\LARGE\textcolor{mainBlue}{#1}}\hfill
{\headfont\large\textcolor{accentColor}{\SubjectName}}\par
\vspace{2pt}
{\large #2}\hfill{\small\textcolor{mainBlue}{厚胤塾}}\par
\vspace{8pt}
\hfill 氏名\ \uline{\hspace{12zw}}\hspace{1zw}%
得点\ \fbox{\rule{0pt}{1.6zh}\hspace{5zw}}
\end{tcolorbox}
\vspace{0.5\Cvs}}

%% ============================================================
%%  定理環境（問題・解答系）
%% ============================================================
\newtheorem{problem}{問題}
\newtheorem{reidai}{例題}
\newtheorem{ans}{答え}
\newtheorem{ensyu}{演習}
\newtheorem{solution}{解答}
\newtheorem{explanation}{解説}
\newtheorem{kakunin}{確認問題}
\newtheorem{kakusolu}{解説}

%% ============================================================
%%  アンダーライン
%% ============================================================
\font\sixly=lasy6
\newcommand{\waveunderline}[1]{%
  \bgroup\markoverwith{\lower3.5pt\hbox{\sixly \char58}}\ULon{#1}}
\newcommand{\dashunderline}[1]{%
  \bgroup\markoverwith{\hbox to .5em{\hrulefill}}\ULon{#1}}
\newcommand{\solidunderline}[1]{\uline{#1}}

%% ============================================================
%%  穴埋め \anaume 系
%% ============================================================
\setlength{\fboxsep}{1pt}
%% 太枠（初出の解答欄）
\newcommand{\anaume}[1]{\setlength{\fboxrule}{1pt}%
  \framebox[40pt][c]{\mbox{\textbf{\large{#1}}}}\setlength{\fboxrule}{0.4pt}}
%% 細枠（2回目以降）
\newcommand{\sanaume}[1]{\framebox[40pt][c]{\mbox{\large{#1}}}}
%% 横長太枠（3桁以上）
\newcommand{\banaume}[1]{\setlength{\fboxrule}{1pt}%
  \fbox{\mbox{\textbf{\large{\phantom{a}#1\phantom{a}}}}}\setlength{\fboxrule}{0.4pt}}
%% 横長細枠
\newcommand{\sbanaume}[1]{\fbox{\mbox{\large{\phantom{a}#1\phantom{a}}}}}
%% 指数用
\newcommand{\sisuuanaume}[1]{\setlength{\fboxrule}{1pt}%
  \framebox[16pt][c]{\mbox{\textbf{#1}}}\setlength{\fboxrule}{0.4pt}}
\newcommand{\ssisuuanaume}[1]{\framebox[16pt][c]{\mbox{#1}}}
%% 選択肢用（二重枠）
\newcommand{\senntakuanaume}[1]{\setlength{\fboxrule}{0.6pt}%
  \doublebox{\mbox{\textbf{\phantom{ア}#1\phantom{ア}}}}\setlength{\fboxrule}{0.4pt}}
\newcommand{\ssenntakuanaume}[1]{\doublebox{\mbox{\phantom{ア}#1\phantom{ア}}}}

%% ============================================================
%%  hyperref（最後に読み込む）
%% ============================================================
\usepackage{hyperref}
\usepackage{pxjahyper}
\hypersetup{hidelinks}
 の直後に %% ============================================================
%% koushin-sekaishi.tex --- 厚胤塾 世界史用プリアンブル
%% 使い方: %% ============================================================
%% koushin-base.tex --- 厚胤塾 共通プリアンブル（全科目共通）
%% ------------------------------------------------------------
%% 対応環境: platex + dvipdfmx（Cloud LaTeX 標準構成）
%%
%% 使い方:
%%   \documentclass[dvipdfmx]{jsbook}   % jsarticle でも可
%%   \input{koushin-base}      % ← 必ず最初に読み込む
%%   \input{koushin-physics}   % ← 科目別プリアンブル（1つだけ）
%%   \begin{document}
%%   ...
%%
%% ブランドカラー: mainBlue = RGB(0,51,102)  ← 全科目共通・不変
%% 科目色       : accentColor ← 科目別ファイルが上書きする
%% ============================================================

%% ---------- 基本パッケージ ----------
\usepackage{graphicx}
\usepackage{wrapfig}
\usepackage{xcolor}
\usepackage{amsmath,amsthm,amssymb}
\usepackage{bm}
\usepackage{enumitem}
\usepackage{fancybox}
\usepackage{multicol}
\usepackage{pdfpages}
\usepackage{float}
\usepackage{marginnote}
\usepackage{ascmac}
\usepackage{url}
\usepackage{tikz}
\usetikzlibrary{patterns}
\usepackage[most]{tcolorbox}
\tcbuselibrary{skins,breakable,xparse}
\usepackage[normalem]{ulem}
\usepackage{tocloft}
\usepackage{fancyhdr}

%% ---------- ブランドカラー ----------
\definecolor{mainBlue}{RGB}{0,51,102}      % 厚胤塾ブランド色（固定）
\definecolor{accentColor}{RGB}{0,51,102}   % 科目色（科目別ファイルで上書き）
\providecommand{\SubjectName}{}            % 科目名（科目別ファイルで上書き）

%% ---------- 余白 ----------
\usepackage[top=25mm,bottom=25mm,left=25mm,right=25mm]{geometry}

%% ---------- 2段組の設定 ----------
\setlength{\columnsep}{6mm}
\setlength{\columnseprule}{0.4pt}
\def\columnseprulecolor{\color{mainBlue}}

%% ---------- 丸数字 \maru{}（platexでUnicode①が使えないため） ----------
\DeclareRobustCommand{\maru}[1]{%
  \tikz[baseline=(char.base)]{%
    \node[shape=circle,draw,inner sep=0.6pt,line width=0.4pt] (char)
      {\footnotesize #1};}}

%% ---------- 式番号・図表番号（丸数字＆通し番号） ----------
\makeatletter
\@ifundefined{c@chapter}{}{%
  \counterwithout{equation}{chapter}%
  \counterwithout{figure}{chapter}%
  \counterwithout{table}{chapter}%
}
\makeatother
\renewcommand{\theequation}{\maru{\arabic{equation}}}
\renewcommand{\thefigure}{\arabic{figure}}
\renewcommand{\thetable}{\arabic{table}}

%% ---------- enumerate のラベル形式 ----------
\setlist[enumerate,1]{label=(\arabic*),  font=\normalfont\upshape}
\setlist[enumerate,2]{label=(\roman*).,  font=\normalfont\upshape}
\setlist[enumerate,3]{label=(\alph*),    font=\normalfont\upshape}

%% ---------- 目次デザイン ----------
\renewcommand{\contentsname}{内容一覧}
\setcounter{tocdepth}{2}
\renewcommand{\cftpartfont}{\LARGE\bfseries\color{mainBlue}}
\renewcommand{\cftpartpagefont}{\LARGE\bfseries\color{mainBlue}}
\renewcommand{\cftsecfont}{\Large\bfseries\color{accentColor}}
\renewcommand{\cftsecpagefont}{\Large\bfseries\color{accentColor}}
\setlength{\cftbeforesecskip}{6pt}

%% ---------- ヘッダー・フッター ----------
\pagestyle{fancy}
\fancyhf{}
\fancyhead[LE]{\textcolor{accentColor}{\textsf{\leftmark}}}
\fancyhead[RO]{\textcolor{accentColor}{\textsf{\rightmark}}}
\fancyfoot[C]{\textcolor{mainBlue}{\thepage}}
\fancyfoot[LE,RO]{\textcolor{accentColor}{\scriptsize 厚胤塾\ \SubjectName}}
\renewcommand{\headrulewidth}{0pt}
\renewcommand{\footrulewidth}{0pt}
\renewcommand{\headrule}{%
  \hbox to \headwidth{\color{accentColor}\leaders\hrule height 0.5pt\hfill}}
\renewcommand{\footrule}{%
  \hbox to \headwidth{\color{accentColor}\leaders\hrule height 0.5pt\hfill}}

%% ============================================================
%%  装飾枠（tcolorbox）--- 枠色は accentColor に連動
%% ============================================================

%% 授業例・例示用（白背景・影つき・タイトルあり）
\newtcolorbox{jyurei}[1][]{
  enhanced, colframe=accentColor, colback=white,
  coltitle=white, fonttitle=\bfseries, drop fuzzy shadow,
  title=#1, boxrule=1pt, arc=4pt, breakable
}

%% グレー枠・タイトル浮き出し
\newtcolorbox{mybox}[1][]{
  enhanced, colframe=white!35!black, colback=white,
  title=#1, fonttitle=\bfseries, breakable=true,
  coltitle=white!20!black,
  attach boxed title to top left={xshift=5mm, yshift=-3mm},
  boxed title style={colframe=white!35!black, colback=white},
  top=4mm
}

%% 科目色枠・クリーム背景・タイトルボックス浮き出し
\newtcolorbox{b_s_b_t}[1][]{
  colback=yellow!10, colframe=accentColor, coltitle=white,
  fonttitle=\bfseries, title=#1, boxrule=1pt, arc=4pt,
  breakable, enhanced,
  attach boxed title to top left={xshift=1em,
    yshift*=-\tcboxedtitleheight/2},
  colbacktitle=accentColor,
  boxed title style={frame code={%
    \path[fill=accentColor] (frame.south west)--(frame.north west)--
    (frame.north east)--(frame.south east)--cycle;},
    interior hidden},
}

%% 科目色実線枠・クリーム背景・シンプルタイトル
\newtcolorbox{bluebox}[1][]{
  colframe=accentColor, colback=yellow!10, coltitle=white,
  fonttitle=\bfseries, title=#1, boxrule=1pt, arc=4pt, breakable
}

%% 破線枠・クリーム背景・タイトルあり
\newenvironment{mybox_d_b}[1]{%
  \begin{tcolorbox}[
    coltitle=black, title=#1, colback=yellow!10,
    frame hidden, interior hidden, enhanced, arc=1em,
    borderline={1pt}{0pt}{dashed},
    attach boxed title to top left={xshift=1em,
      yshift*=-\tcboxedtitleheight/2},
    colbacktitle=white, boxed title style={frame hidden},
    breakable,
  ]
}{\end{tcolorbox}}

%% 破線枠・タイトルなし・背景なし
\newenvironment{mybox_dash}{%
  \begin{tcolorbox}[
    frame hidden, interior hidden, enhanced, arc=1em,
    borderline={1pt}{0pt}{dashed}
  ]
}{\end{tcolorbox}}

%% 破線枠・タイトルあり・背景なし（解法など）
\newenvironment{kaiho}[1]{%
  \begin{tcolorbox}[
    coltitle=black, title=#1, frame hidden, interior hidden,
    enhanced, arc=1em, borderline={1pt}{0pt}{dashed},
    attach boxed title to top left={xshift=1em,
      yshift*=-\tcboxedtitleheight/2},
    colbacktitle=white, boxed title style={frame hidden},
  ]
}{\end{tcolorbox}}

%% 黒実線・タイトル切り欠き枠（原理・法則）
\DeclareTColorBox{genri}{o m O{.5} O{}}{%
  empty, left=4mm, right=4mm, top=-1mm,
  attach boxed title to top left={xshift=1.2zw},
  boxed title style={empty, left=-2mm, right=-2mm},
  colframe=black, coltitle=black, coltext=black, breakable,
  underlay unbroken={%
    \draw[black, line width=#3pt]
      (title.east) -- (title.east-|frame.east) --
      (frame.south east) -- (frame.south west) --
      (title.west-|frame.west) -- (title.west);},
  underlay first={%
    \draw[black, line width=#3pt]
      (title.east) -- (title.east-|frame.east) -- (frame.south east);
    \draw[black, line width=#3pt]
      (frame.south west) -- (title.west-|frame.west) -- (title.west);},
  underlay middle={%
    \draw[black, line width=#3pt] (frame.north east) -- (frame.south east);
    \draw[black, line width=#3pt] (frame.south west) -- (frame.north west);},
  underlay last={%
    \draw[black, line width=#3pt]
      (frame.north east) -- (frame.south east) --
      (frame.south west) -- (frame.north west);},
  fonttitle=\gtfamily,
  IfValueTF={#1}{title=【#2】〈#1〉}{title=【#2】},
  #4
}

%% 警告枠（赤：全科目共通で赤固定）
\newtcolorbox{keikoku}[1][]{
  colframe=red!70!black, colback=yellow!10, coltitle=white,
  fonttitle=\bfseries, title=#1, boxrule=1pt, arc=4pt,
  breakable, enhanced,
  attach boxed title to top left={xshift=1em,
    yshift*=-\tcboxedtitleheight/2},
  colbacktitle=red!70!black,
  boxed title style={frame code={%
    \path[fill=red!70!black] (frame.south west)--(frame.north west)--
    (frame.north east)--(frame.south east)--cycle;},
    interior hidden},
}

%% 要点まとめ枠（科目色・薄い科目色背景）
\newtcolorbox{pointbox}[1][ポイント]{
  enhanced, colframe=accentColor, colback=accentColor!4!white,
  coltitle=white, colbacktitle=accentColor, fonttitle=\bfseries,
  title=#1, boxrule=1pt, arc=3pt, breakable,
  attach boxed title to top left={xshift=1em,
    yshift*=-\tcboxedtitleheight/2}
}

%% ============================================================
%%  見出しデザイン
%% ============================================================

%% --- セクション（番号白抜き＋帯） ---
\makeatletter
\renewcommand{\thesection}{\arabic{section}}
\newcommand\Section[1]{%
  \Section@topmargin%
  \refstepcounter{section}%
  {\setlength{\parindent}{0pt}%
   \Section@font%
   \Section@title{\thesection}{#1}%
   \par\nobreak}%
  \Section@bottommargin%
}
\newcommand{\Section@topmargin}{\vspace{2\Cvs}}
\newcommand{\Section@bottommargin}{\vspace{1\Cvs}}
\newcommand{\Section@font}{\headfont\Large}
\newcommand{\Section@title}[2]{%
  \textcolor{accentColor!30}{\rlap{\rule[-0.6zh]{\textwidth}{2zh}}}%
  \textcolor{accentColor}{\rlap{\rule[-0.6zh]{3.5zw}{2zh}}}%
  \makebox[3.5zw][c]{\textcolor{white}{#1}}%
  \hspace{0.9zw}#2%
}
\makeatother

%% --- 太バー付き見出し ---
\newcommand{\DecoratedSection}[1]{%
\vspace{2\Cvs}%
{\setlength{\parindent}{0pt}%
\begin{tcolorbox}[
  enhanced, colframe=accentColor, colback=white,
  boxrule=1pt, arc=0mm, left=10pt, right=10pt, top=8pt, bottom=8pt,
  leftrule=12pt, width=\textwidth, fontupper=\headfont\Large
]{#1}\end{tcolorbox}}%
\vspace{\Cvs}}

%% --- 二重ライン見出し ---
\newcommand{\FancySection}[1]{%
\vspace{2\Cvs}%
{\setlength{\parindent}{0pt}%
\begin{tcolorbox}[
  enhanced, colframe=accentColor, colback=white,
  boxrule=1pt, arc=3mm, left=10pt, right=10pt, top=8pt, bottom=8pt,
  leftrule=0pt, width=\textwidth,
  overlay={%
    \draw[accentColor, line width=3pt]
      ([xshift=2pt]frame.north west) -- ([xshift=-2pt]frame.north east);
    \draw[accentColor!30, line width=1.5pt]
      ([xshift=2pt, yshift=-3pt]frame.north west) --
      ([xshift=-2pt, yshift=-3pt]frame.north east);},
  fontupper=\headfont\Large
]{#1}\end{tcolorbox}}%
\vspace{\Cvs}}

%% --- 入試問題バナー ---
\newcommand{\ChallengeTitle}{%
\vspace{0.3cm}
\begin{tcolorbox}[
  enhanced, colframe=accentColor, colback=white,
  boxrule=0pt, arc=0mm, left=20pt, right=20pt, top=8pt, bottom=8pt,
  width=\textwidth,
  overlay={%
    \fill[accentColor]
      ([xshift=-3pt]frame.north west) rectangle (frame.south west);
    \fill[accentColor!90]
      (frame.north west) rectangle ([xshift=\textwidth]frame.north east);
    \fill[left color=accentColor!20, right color=white]
      ([yshift=-3pt]frame.north west) rectangle
      ([xshift=8cm, yshift=-4pt]frame.north west);
    \node[anchor=west, inner sep=0pt, font=\Large\bfseries]
      at ([xshift=25pt, yshift=-14pt]frame.north west)
      {\textcolor{accentColor}{入試問題に挑戦}};},
  height=2cm
]\end{tcolorbox}
\vspace{0.3cm}}

%% --- テスト用タイトルバナー（#1: テスト名, #2: 範囲） ---
\newcommand{\KoushinTestTitle}[2]{%
\begin{tcolorbox}[
  enhanced, colframe=mainBlue, colback=white,
  boxrule=1pt, arc=0mm, left=16pt, right=12pt, top=10pt, bottom=10pt,
  overlay={%
    \fill[accentColor]
      (frame.north west) rectangle ([xshift=6pt]frame.south west);}
]
{\headfont\LARGE\textcolor{mainBlue}{#1}}\hfill
{\headfont\large\textcolor{accentColor}{\SubjectName}}\par
\vspace{2pt}
{\large #2}\hfill{\small\textcolor{mainBlue}{厚胤塾}}\par
\vspace{8pt}
\hfill 氏名\ \uline{\hspace{12zw}}\hspace{1zw}%
得点\ \fbox{\rule{0pt}{1.6zh}\hspace{5zw}}
\end{tcolorbox}
\vspace{0.5\Cvs}}

%% ============================================================
%%  定理環境（問題・解答系）
%% ============================================================
\newtheorem{problem}{問題}
\newtheorem{reidai}{例題}
\newtheorem{ans}{答え}
\newtheorem{ensyu}{演習}
\newtheorem{solution}{解答}
\newtheorem{explanation}{解説}
\newtheorem{kakunin}{確認問題}
\newtheorem{kakusolu}{解説}

%% ============================================================
%%  アンダーライン
%% ============================================================
\font\sixly=lasy6
\newcommand{\waveunderline}[1]{%
  \bgroup\markoverwith{\lower3.5pt\hbox{\sixly \char58}}\ULon{#1}}
\newcommand{\dashunderline}[1]{%
  \bgroup\markoverwith{\hbox to .5em{\hrulefill}}\ULon{#1}}
\newcommand{\solidunderline}[1]{\uline{#1}}

%% ============================================================
%%  穴埋め \anaume 系
%% ============================================================
\setlength{\fboxsep}{1pt}
%% 太枠（初出の解答欄）
\newcommand{\anaume}[1]{\setlength{\fboxrule}{1pt}%
  \framebox[40pt][c]{\mbox{\textbf{\large{#1}}}}\setlength{\fboxrule}{0.4pt}}
%% 細枠（2回目以降）
\newcommand{\sanaume}[1]{\framebox[40pt][c]{\mbox{\large{#1}}}}
%% 横長太枠（3桁以上）
\newcommand{\banaume}[1]{\setlength{\fboxrule}{1pt}%
  \fbox{\mbox{\textbf{\large{\phantom{a}#1\phantom{a}}}}}\setlength{\fboxrule}{0.4pt}}
%% 横長細枠
\newcommand{\sbanaume}[1]{\fbox{\mbox{\large{\phantom{a}#1\phantom{a}}}}}
%% 指数用
\newcommand{\sisuuanaume}[1]{\setlength{\fboxrule}{1pt}%
  \framebox[16pt][c]{\mbox{\textbf{#1}}}\setlength{\fboxrule}{0.4pt}}
\newcommand{\ssisuuanaume}[1]{\framebox[16pt][c]{\mbox{#1}}}
%% 選択肢用（二重枠）
\newcommand{\senntakuanaume}[1]{\setlength{\fboxrule}{0.6pt}%
  \doublebox{\mbox{\textbf{\phantom{ア}#1\phantom{ア}}}}\setlength{\fboxrule}{0.4pt}}
\newcommand{\ssenntakuanaume}[1]{\doublebox{\mbox{\phantom{ア}#1\phantom{ア}}}}

%% ============================================================
%%  hyperref（最後に読み込む）
%% ============================================================
\usepackage{hyperref}
\usepackage{pxjahyper}
\hypersetup{hidelinks}
 の直後に %% ============================================================
%% koushin-sekaishi.tex --- 厚胤塾 世界史用プリアンブル
%% 使い方: \input{koushin-base} の直後に \input{koushin-sekaishi}
%% （中身は koushin-history.tex 共通、科目名のみ変更）
%% ============================================================
\input{koushin-history}
\renewcommand{\SubjectName}{世界史}

%% （中身は koushin-history.tex 共通、科目名のみ変更）
%% ============================================================
%% ============================================================
%% koushin-history.tex --- 厚胤塾 歴史（日本史・世界史）共通プリアンブル
%% 使い方: 通常はこのファイルを直接 \input せず、
%%   \input{koushin-nihonshi}  または  \input{koushin-sekaishi}
%% を使う（科目名だけ切り替わる）。
%% 科目色: 金茶 RGB(133,94,33)
%% ============================================================

\definecolor{accentColor}{RGB}{133,94,33}
\providecommand{\SubjectName}{}\renewcommand{\SubjectName}{歴史}

%% ---------- ルビ ----------
%% \ruby{漢字}{かん|じ} のように使う
\usepackage{pxrubrica}

%% ---------- 歴史用マクロ ----------
%% 年号強調（例: \nen{1867} → 太字＋科目色）
\newcommand{\nen}[1]{\textcolor{accentColor}{\textbf{#1年}}}

%% 重要語句（例: \juuyou{大政奉還}）
\newcommand{\juuyou}[1]{\textcolor{accentColor}{\textgt{#1}}}

%% ---------- 史料ボックス ----------
\newtcolorbox{shiryo}[1][史料]{
  enhanced, colframe=accentColor, colback=accentColor!4!white,
  coltitle=white, colbacktitle=accentColor, fonttitle=\bfseries,
  title=#1, boxrule=1pt, arc=3pt, breakable,
  attach boxed title to top left={xshift=1em,
    yshift*=-\tcboxedtitleheight/2}
}

%% ---------- 年表ボックス ----------
%% 中で tabular を使う想定:
%%   \begin{nenpyo}[幕末年表]
%%     \begin{tabular}{ll}
%%       1853年 & ペリー来航 \\ ...
%%     \end{tabular}
%%   \end{nenpyo}
\newtcolorbox{nenpyo}[1][年表]{
  enhanced, colframe=accentColor, colback=white,
  coltitle=white, colbacktitle=accentColor, fonttitle=\bfseries,
  title=#1, boxrule=1pt, arc=3pt, breakable,
  attach boxed title to top left={xshift=1em,
    yshift*=-\tcboxedtitleheight/2}
}

\renewcommand{\SubjectName}{世界史}

%% （中身は koushin-history.tex 共通、科目名のみ変更）
%% ============================================================
%% ============================================================
%% koushin-history.tex --- 厚胤塾 歴史（日本史・世界史）共通プリアンブル
%% 使い方: 通常はこのファイルを直接 \input せず、
%%   \input{koushin-nihonshi}  または  %% ============================================================
%% koushin-sekaishi.tex --- 厚胤塾 世界史用プリアンブル
%% 使い方: \input{koushin-base} の直後に \input{koushin-sekaishi}
%% （中身は koushin-history.tex 共通、科目名のみ変更）
%% ============================================================
\input{koushin-history}
\renewcommand{\SubjectName}{世界史}

%% を使う（科目名だけ切り替わる）。
%% 科目色: 金茶 RGB(133,94,33)
%% ============================================================

\definecolor{accentColor}{RGB}{133,94,33}
\providecommand{\SubjectName}{}\renewcommand{\SubjectName}{歴史}

%% ---------- ルビ ----------
%% \ruby{漢字}{かん|じ} のように使う
\usepackage{pxrubrica}

%% ---------- 歴史用マクロ ----------
%% 年号強調（例: \nen{1867} → 太字＋科目色）
\newcommand{\nen}[1]{\textcolor{accentColor}{\textbf{#1年}}}

%% 重要語句（例: \juuyou{大政奉還}）
\newcommand{\juuyou}[1]{\textcolor{accentColor}{\textgt{#1}}}

%% ---------- 史料ボックス ----------
\newtcolorbox{shiryo}[1][史料]{
  enhanced, colframe=accentColor, colback=accentColor!4!white,
  coltitle=white, colbacktitle=accentColor, fonttitle=\bfseries,
  title=#1, boxrule=1pt, arc=3pt, breakable,
  attach boxed title to top left={xshift=1em,
    yshift*=-\tcboxedtitleheight/2}
}

%% ---------- 年表ボックス ----------
%% 中で tabular を使う想定:
%%   \begin{nenpyo}[幕末年表]
%%     \begin{tabular}{ll}
%%       1853年 & ペリー来航 \\ ...
%%     \end{tabular}
%%   \end{nenpyo}
\newtcolorbox{nenpyo}[1][年表]{
  enhanced, colframe=accentColor, colback=white,
  coltitle=white, colbacktitle=accentColor, fonttitle=\bfseries,
  title=#1, boxrule=1pt, arc=3pt, breakable,
  attach boxed title to top left={xshift=1em,
    yshift*=-\tcboxedtitleheight/2}
}

\renewcommand{\SubjectName}{世界史}

%% を使う（科目名だけ切り替わる）。
%% 科目色: 金茶 RGB(133,94,33)
%% ============================================================

\definecolor{accentColor}{RGB}{133,94,33}
\providecommand{\SubjectName}{}\renewcommand{\SubjectName}{歴史}

%% ---------- ルビ ----------
%% \ruby{漢字}{かん|じ} のように使う
\usepackage{pxrubrica}

%% ---------- 歴史用マクロ ----------
%% 年号強調（例: \nen{1867} → 太字＋科目色）
\newcommand{\nen}[1]{\textcolor{accentColor}{\textbf{#1年}}}

%% 重要語句（例: \juuyou{大政奉還}）
\newcommand{\juuyou}[1]{\textcolor{accentColor}{\textgt{#1}}}

%% ---------- 史料ボックス ----------
\newtcolorbox{shiryo}[1][史料]{
  enhanced, colframe=accentColor, colback=accentColor!4!white,
  coltitle=white, colbacktitle=accentColor, fonttitle=\bfseries,
  title=#1, boxrule=1pt, arc=3pt, breakable,
  attach boxed title to top left={xshift=1em,
    yshift*=-\tcboxedtitleheight/2}
}

%% ---------- 年表ボックス ----------
%% 中で tabular を使う想定:
%%   \begin{nenpyo}[幕末年表]
%%     \begin{tabular}{ll}
%%       1853年 & ペリー来航 \\ ...
%%     \end{tabular}
%%   \end{nenpyo}
\newtcolorbox{nenpyo}[1][年表]{
  enhanced, colframe=accentColor, colback=white,
  coltitle=white, colbacktitle=accentColor, fonttitle=\bfseries,
  title=#1, boxrule=1pt, arc=3pt, breakable,
  attach boxed title to top left={xshift=1em,
    yshift*=-\tcboxedtitleheight/2}
}
